% 03_method.tex
\vspace{-.2cm}
\section{Method}
\label{sec:method}

For a speech waveform $\mathbf{x}$, a frozen no-reference predictor $f$ returns the scalar MOS estimate $\qhat=f(\mathbf{x})$. We extract a fixed-dimensional representation $\zvec=\phi(\mathbf{x})\in\mathbb{R}^{d}$ from each predictor's penultimate layer, mean-pooling frame- or segment-level activations: the second dense layer before the final head for DNSMOS ($d=64$), the self-attention time-dependency module before the pooling heads for NISQA ($d=64$), and the BiLSTM output of the LDConditioner for UTMOS ($d=1024$). Each predictor has its own feature space and therefore its own reference distribution.

\vspace{-.2cm}
\subsection{Feature-Space Scores}
\label{ssec:features}

For each predictor, T-SIM supplies reference features $\{\zvec_i\}_{i=1}^{N}$, where $N=10{,}000$. The maximum-likelihood empirical mean and covariance are
\begin{align}
    \hat{\muvec}
    &= \frac{1}{N}\sum_{i=1}^{N}\zvec_i,
    \label{eq:mean} \\
    \hat{\Sigmat}
    &= \frac{1}{N}\sum_{i=1}^{N}
    (\zvec_i-\hat{\muvec})(\zvec_i-\hat{\muvec})^\top.
    \label{eq:cov}
\end{align}
Here $N>d$ for all predictors, with $d=64$ for DNSMOS and NISQA ($N/d=156$) and $d=1024$ for UTMOS ($N/d=9.8$). We use the empirical covariance throughout and do not apply Ledoit-Wolf shrinkage~\cite{LedoitWolf2004}. At UTMOS's sample size, shrinkage changes AUROC by at most 0.09 (domain) and 0.12 (error), both above the 0.075 spread across scores in Table~\ref{tab:ood-score}, so per-predictor rankings are not resolvable here. The near-chance conclusion of Section~\ref{ssec:detection-results} is unchanged under either estimator. The Mahalanobis score is
\begin{align}
    d_M(\mathbf{x})
    =
    \left[
        (\zvec-\hat{\muvec})^\top
        \hat{\Sigmat}^{-1}
        (\zvec-\hat{\muvec})
    \right]^{1/2}.
    \label{eq:mahal}
\end{align}
We also evaluate Euclidean distance $\lVert\zvec-\hat{\muvec}\rVert_2$ and the distance to the nearest T-SIM representation~\cite{Sun2022}, fixing $k=1$. A repo-only sweep over $k\in\{1,\ldots,50\}$ changes mean AUROC by at most 0.004. A larger score is prespecified to mean farther from the reference distribution. These scores quantify representation shift rather than estimate MOS error.

\vspace{-.2cm}
\subsection{Split Conformal Interval}
\label{ssec:conformal}

V-SIM is partitioned by content identifier into disjoint calibration and in-domain evaluation halves of 1{,}250 clips each. Let
$\mathcal{D}_{\mathrm{cal}}=\{(\mathbf{x}_j,q_j^*)\}_{j=1}^{m}$ denote the calibration half, where $m=1{,}250$ and $q_j^*$ is subjective MOS. The absolute-residual nonconformity score is
\begin{align}
    s_j = \lvert q_j^*-\hat{q}_j\rvert,
    \qquad \hat{q}_j=f(\mathbf{x}_j).
    \label{eq:nonconf}
\end{align}
Let $s_{(1)}\leq\cdots\leq s_{(m)}$ be the ordered calibration scores and let $\alpha\in(0,1)$ denote the miscoverage rate, so nominal coverage is $1-\alpha$ (we use $\alpha=0.10$). For $\alpha\geq 1/(m+1)$, define
\begin{align}
    k_\alpha
    &= \left\lceil(m+1)(1-\alpha)\right\rceil,% \\
    &\hat Q_{1-\alpha} = s_{(k_\alpha)}.
    \label{eq:conformal-quantile}
\end{align}
The fixed-width interval is
\begin{align}
    \Calpha(\mathbf{x})
    =
    [\qhat-\hat Q_{1-\alpha},\,
     \qhat+\hat Q_{1-\alpha}].
    \label{eq:interval}
\end{align}
We do not clip the endpoints to $[1,5]$, since clipping can alter coverage at the scale boundaries. If the calibration examples and a new example are exchangeable and calibration labels were not used to fit or select $f$, then split conformal prediction gives the marginal guarantee
\begin{align}
    \Pr\!\left\{\qtrue\in\Calpha(\mathbf{x})\right\}
    \geq 1-\alpha.
    \label{eq:coverage-guarantee}
\end{align}
The content-disjoint split prevents direct content leakage but does not itself prove exchangeability, which can in principle be tested online~\cite{Vovk2003mondrian}, though we do not conduct that test here. The NISQA model was developed using NISQA Corpus partitions~\cite{Mittag2021nisqa}, so its intervals are treated as empirical rather than as a fresh-data guarantee.

\vspace{-.2cm}
\subsection{Distance-Stratified Interval}
\label{ssec:adaptive}

Following Mondrian conformal regression with binned difficulty estimates~\cite{Bostrom2020Mondrian}, we test whether Mahalanobis distance supports instance-dependent widths. We select Mahalanobis distance because Section~\ref{ssec:detection-results} shows it is comparable to nearest-neighbor distance on domain-detection AUROC while needing only a fixed $(d,d)$ covariance matrix rather than a reference bank that grows with the calibration set. The calibration examples are divided into $K$ empirical quantile bins of $d_M(\mathbf{x})$. Let $I_k$ be the calibration indices assigned to bin $k$, let $m_k=|I_k|$, and let $s_{(r)}^{(k)}$ be the $r$th ordered residual within that bin. For $\alpha\geq1/(m_k+1)$,
\begin{align}
    \ell_{\alpha,k}
    &= \left\lceil(m_k+1)(1-\alpha)\right\rceil, %\\
    &\hat Q_{1-\alpha}^{(k)}=s_{(\ell_{\alpha,k})}^{(k)}.
    \label{eq:bin-quantile}
\end{align} 
If $b(\mathbf{x})\in\{1,\ldots,K\}$ is the bin assigned using the calibration-derived boundaries, then
\begin{align}
    \Calpha^{\mathrm{adapt}}(\mathbf{x})
    =
    [\qhat-\hat Q_{1-\alpha}^{(b(\mathbf{x}))},\,
     \qhat+\hat Q_{1-\alpha}^{(b(\mathbf{x}))}].
    \label{eq:adaptive}
\end{align}
Scores outside the calibration range enter the first or last bin. We use $K=5$ and hold the boundaries and radii fixed at evaluation. Because the boundaries are estimated from the calibration covariates and shifted test examples need not be exchangeable with calibration examples, we treat Eq.~\eqref{eq:adaptive} as a diagnostic post-hoc construction. We do not claim conditional or arbitrary-shift coverage for it.
